\documentclass[10pt,a4paper]{amsart}
\usepackage[margin=19mm]{geometry}
\usepackage{amsmath,amssymb,mathtools}
\usepackage[scr=rsfs,cal=euler]{mathalfa}
\usepackage{xfrac}
\usepackage[unicode,hidelinks,bookmarks=false]{hyperref}
\newcommand{\ie}{i.e.\ }
\newcommand{\cf}{cf.\ }
\newcommand{\resp}{respectively\ }
\newcommand{\Subsection}{\subsection}
\providecommand{\nobreakdash}{\nobreak\hbox{-}}
\setlength{\emergencystretch}{2em}
\multlinegap0pt
\usepackage{amssymb}
\usepackage{mathtools}
\usepackage{xcolor}
\newtheorem{prop}{Proposition}
\title{Multisets with few special directions and small weight codewords in Desarguesian planes}
\author{Sam Adriaensen, Tam\'as Sz\H{o}nyi, Zsuzsa Weiner}
\date{Source: arXiv:2411.19201v3 (CC BY 4.0)}
\begin{document}
\maketitle

\noindent\emph{Source and licence.} Multisets with few special directions and small weight codewords in Desarguesian planes. Version arXiv:2411.19201v3 (CC BY 4.0), distributed by arXiv under the Creative Commons Attribution 4.0 International licence, http://creativecommons.org/licenses/by/4.0/, as recorded in the arXiv licence metadata for this exact version and shown on the abstract page https://arxiv.org/abs/2411.19201v3. Full licence notice: \texttt{licenses/arxiv-2411.19201-CC-BY-4.0.txt}.\par\smallskip

\noindent\textit{Editorial scope.} Counted unit: the article's prop 2\_spec\_dirs (source0.tex lines 658--661). The article's complete original proof of it is delivered with it (source0.tex lines 663--687, one \texttt{proof} environment of 25 lines). No mathematics is written, repaired or re-derived: the statement and the proof are the article's own lines, in the article's own notation, and no retained line is substituted or re-flowed.  No further source-local context is needed: every cross-reference of every retained line resolves inside the two delivered spans.  Nothing was omitted from any retained span, so the package carries no omission record.  No retained line contains a citation, so the wrapper carries no bibliography and no bibliography entry.  No retained line contains a figure environment, an image-inclusion command or drawing code, so no image is delivered and the package declares no asset dependency.  Editorial formatting only, in addition: an AMS \texttt{amsart} preamble that restates the article's own declarations and macros used by the retained lines, namely the environments \texttt{prop} on separate counters, together with the standard packages those lines need.  The retained environments print in this document as Proposition 1 in order of appearance, whereas the article prints the same items with the numbering of its own full document.  The complete native source of the article as distributed in the arXiv e-print for this exact version is bundled unchanged as \texttt{sources/169292/source0.tex}.
\begin{prop}
\label{2_spec_dirs}
 Suppose that a multiset $M$ of size $nq$ determines exactly $2$ special directions. Then $M$ is the union of lines through these special directions.
\end{prop}

\begin{proof}
 We may suppose that the special directions are $(0)$ and $(\infty)$.
 Let $\mu(x,y)$ denote the multiplicity of $M$ at the point $(x,y)$.
 Suppose that the lines $X = x$ and $Y = y$ intersect $M$ in $a_x$ and $b_y$ points, respectively.
 By counting the lines through a point $(x,y)$, we see that
 \[
  (q-1) n + a_x + b_y = qn + q \mu(x,y).
 \]
 In other words,
 \[
  \mu(x,y) = \frac{a_x + b_y - n}{q}.
 \]
 This means that $a_x + b_y \equiv n \pmod{q}$ for all $x$ and $y$.
 Fixing $y$ and varying $x$, we see that there exist integers $r_x$ such that $a_x = a + q r_x$ for some $a$ with $0 \leq a < q$.
 Likewise, $b_y = b + q s_y$ for some integers $s_y$ and some $b$ with $0 \leq b < q$.
 Moreover, $a + b \equiv n \pmod q$, hence $a+b = n - qt$ for some integer $t$.

 Plugging this into the equation, we see that
 \[
  \mu(x,y) = r_x + s_y - t.
 \]
Take $x$ with $r_x$ minimal and $y$ with $s_y$ minimal.
Then $\mu(x,y) \geq 0$, hence we can write $t = t_1 + t_2$ with $0 \leq t_1 \leq r_x$ for each $x$ and $0 \leq t_2 \leq r_y$ for each $y$.
This means that we can obtain $M$ as the union of $r_x - t_1$ times the line $X = x$ and $s_y - t_2$ times the line $Y = y$ for all $x$ and $y$.
\end{proof}

\end{document}
